\section{\scshape Introduction}\label{sec:introduction}

\subsection{Context}
\begin{frame}{Context}
	\begin{itemize}
		\item 3D object recognition is a challenging task that may require active perception of the environment for gathering additional sensor information when the level of confidence in the perception analysis is not enough
		\item Simulating sensor 3D depth data from a set of representative environments can help decide the type, number and disposition of sensors that maximize the target objects observable surface area
	\end{itemize}
\end{frame}


\subsection{Research Areas}
\begin{frame}{Research Areas}
	\begin{itemize}
		\item 3D modeling and simulation of:
		\begin{itemize}
			\item The 3D environment, including physics
			\item The different types of depth sensors
		\end{itemize}
		\item 3D rendering with occlusions
		\item Active perception
		\item Object recognition
		\item Bin picking
	\end{itemize}
\end{frame}

\subsection{Software dependencies}
\begin{frame}{Software dependencies}
	\begin{itemize}
		\item Robot Operating System (ROS)
		\begin{itemize}
			\item For fast integration between simulation and real robots and 3D visual inspection (Rviz)
		\end{itemize}
		\item Gazebo simulator
		\begin{itemize}
			\item For world / sensor simulation and 3D rendering
		\end{itemize}
		\item Point Cloud Library (PCL)
		\begin{itemize}
			\item For point cloud processing
		\end{itemize}
	\end{itemize}
	\begin{figure}[!ht]
		\centering
		\includegraphics[height=.08\textheight]{logos/ros-logo}
		\hspace{1.5em}
		\includegraphics[height=.08\textheight]{logos/gazebo-logo}
		\hspace{1.5em}
		\includegraphics[height=.08\textheight]{logos/pcl-logo}
		\caption{Main software dependencies}
	\end{figure}
\end{frame}
